% Generated draft: requires mathematical and source review.
\documentclass[11pt]{article}
\usepackage[margin=25mm]{geometry}
\usepackage{fontspec}
\setmainfont{Latin Modern Roman}
\usepackage{amsmath,amssymb,mathtools,unicode-math}
\setmathfont{Latin Modern Math}
\usepackage{enumitem,needspace,xurl,hyperref}
\setlength{\parindent}{0pt}
\setlength{\parskip}{5pt}
\setlength{\emergencystretch}{3em}
\newcommand{\N}{\mathbb{N}}
\newcommand{\Z}{\mathbb{Z}}
\newcommand{\Q}{\mathbb{Q}}
\newcommand{\R}{\mathbb{R}}
\newcommand{\C}{\mathbb{C}}
\newcommand{\F}{\mathbb{F}}
\newcommand{\A}{\mathbb{A}}
\newcommand{\PP}{\mathbb{P}}
\newcommand{\E}{\mathbb{E}}
\newcommand{\eps}{\varepsilon}
\newcommand{\dd}{\,\mathrm{d}}
\begin{document}
\section*{UnsolvedMath ID: 30000857}
\textbf{Source problem code:} catalogue entry
\subsection{Definitions and mathematical statement}
Let $M=\Gamma\backslash G$ be a compact nilmanifold, where $G$ is a nilpotent real Lie group and $\Gamma$ a discrete cocompact subgroup, equipped with a left-invariant integrable complex structure $J$. Such a structure is induced by an endomorphism of the Lie algebra with $J^2=-1$ and vanishing Nijenhuis tensor. The usual connected/simply-connected convention on $G$ is not explicit in the source and should be confirmed.

The source asks three tasks:
\begin{enumerate}
\item Determine the small deformations of the complex manifold $(M,J)$. This concerns its full local complex-deformation space, not merely the subfamily of left-invariant complex structures. In particular, criteria ensuring that all small deformations remain left invariant are only part of the requested picture.
\item Determine when $(M,J)$ has a geometric description as an iterated principal holomorphic torus bundle. In the usual sense this is a finite tower of holomorphic principal bundles with compact complex torus fibers, ending at a point or the equivalent terminal-torus convention. Principal means translation structure group; an arbitrary holomorphic torus fibration is not silently identified with a principal bundle.
\item Determine all deformations in the large of such bundles. This asks for the global deformation class in the intended moduli/deformation sense, not just a neighborhood of one complex structure.
\end{enumerate}

The record does not prescribe a Kuranishi or other local deformation model, marked versus unmarked equivalence, the global deformation relation, the allowed torus-bundle tower, or hypotheses involving the lattice and complex Lie-algebra data. Nor does it assert that every deformation remains invariant, nilmanifold-like, or a principal torus bundle. Those are potential conclusions to analyze, not assumptions to add.

The supplied assessment attributes criteria for small-deformation invariance and substantial large-deformation results to Rollenske, with a fairly complete complex-dimension-three picture. The source's general three-part task is not restricted to dimension three or identified with those reported partial theorems.
\subsection{Short English statement}
The entry asks for local deformations of invariant complex nilmanifolds, criteria for an iterated principal holomorphic torus-bundle description, and the full large-deformation class of such bundles. Full complex deformations must remain distinct from the invariant subfamily.
\subsection{Sources}
\begin{itemize}
\item Local UnsolvedMath registry, permanent ID 30000857, code catalogue entry, ``Deformations of Complex Nilmanifolds and Torus Bundles.'' Proposed year 2007; proposer unspecified.
\item \emph{Komplexe Algebraische Geometrie} (2008), Oberwolfach Reports, pp.\ 2791--2850, DOI 10.4171/owr/2007/47; \url{https://doi.org/10.4171/owr/2007/47}. The displayed report year 2008 differs from the DOI and proposed year 2007.
\item Supplied ``Rollenske, nilmanifolds and deformations'': \url{https://doi.org/10.1112/plms/pdp001}.
\item Supplied ``Rollenske, Lie-algebra Dolbeault cohomology'': \url{https://arxiv.org/abs/0709.0467}.
\item Statement status corrected\_verified; the supplied note says official-report content was checked and wording/notation normalized without changing the claim.
\item Literature triage checked 2026-08-21: partially solved. It reports Rollenske's small-deformation invariance criteria, large-deformation results for iterated principal holomorphic torus bundles, and a fairly complete dimension-three picture, while the full classification is broader. Exact hypotheses and current status require independent verification. Created 2026-08-25; updated 2026-09-16.
\item Registry curation by Ulam AI / UnsolvedMath Contributors, CC BY 4.0; underlying report and literature provenance remain as supplied. Unreviewed statement-only draft.
\end{itemize}
\end{document}
